\chapter{Architektura systemu}

\begin{figure}[H]
    \centering
    \makeatletter
    \tikzset{
        database/.style={
            path picture={
                % Usunięto fill – teraz kształt jest przezroczysty, a linie są thick
                \draw[thick] (0, 1.5*\database@segmentheight) circle [x radius=\database@radius,y radius=\database@aspectratio*\database@radius];
                \draw[thick] (-\database@radius, 0.5*\database@segmentheight) arc [start angle=180,end angle=360,x radius=\database@radius, y radius=\database@aspectratio*\database@radius];
                \draw[thick] (-\database@radius,-0.5*\database@segmentheight) arc [start angle=180,end angle=360,x radius=\database@radius, y radius=\database@aspectratio*\database@radius];
                \draw[thick] (-\database@radius,1.5*\database@segmentheight) -- ++(0,-3*\database@segmentheight) arc [start angle=180,end angle=360,x radius=\database@radius, y radius=\database@aspectratio*\database@radius] -- ++(0,3*\database@segmentheight);
            },
            minimum width=2*\database@radius + \pgflinewidth,
            minimum height=3*\database@segmentheight + 2*\database@aspectratio*\database@radius + \pgflinewidth,
        },
        database segment height/.store in=\database@segmentheight,
        database radius/.store in=\database@radius,
        database aspect ratio/.store in=\database@aspectratio,
        database segment height=0.4cm,
        database radius=1.1cm,
        database aspect ratio=0.25,
        service/.style={rectangle, rounded corners=6pt, minimum width=3cm, minimum height=1.5cm, text centered, draw=blue!60!black, fill=blue!5, thick, align=center},
        flow/.style={thick, ->, >=stealth, draw=black!70}
    }
    \makeatother

    \begin{tikzpicture}[node distance=3.0cm and 2.0cm, font=\sffamily]
        
        % --- Definicje komponentów ---
        \node (nginx) [service] {\textbf{Nginx} \\ \small Proxy / Router};
        
        \node (frontend) [service, right=of nginx] {\textbf{Frontend} \\ \small Next.js};
        
        \node (backend) [service, below=of nginx, yshift=-0.3cm] {\textbf{Backend} \\ \small Spring Boot};
        
        \node (detector) [service, left=of backend] {\textbf{Detector} \\ \small Python Services};
        
        % Dodano align=center wewnątrz opcji label
        \node (db) [database, below=of backend, yshift=-0.5cm, label={[align=center, yshift=-0.1cm]below:\textbf{PostgreSQL}}] {};

        % --- Przepływ danych ---
        \draw [flow] (nginx.east) -- (frontend.west);
        \draw [flow] (nginx.south) -- (backend.north);
        \draw [flow] (frontend.south) |- (backend.east);
        \draw [flow] (detector.east) -- (backend.west);
        \draw [flow, <->] (backend.south) -- (db.north);

    \end{tikzpicture}
    \caption{Architektura systemu.}\label{fig:architektura_systemu}
\end{figure}

Przedstawiona na \myref[rysunku]{fig:architektura_systemu} architektura systemu opiera się o następujące kontenery dockera:
\begin{itemize}
\item \textbf{Nginx} --- odgrywa rolę odwróconego proxy, kierując ruch do odpowiednich usług w zależności od ścieżki \texttt{URL}.\ Umożliwia wdrożenie certyfikatu \texttt{SSL} i komunikację poprzez protokół \texttt{HTTPS}.
\item \textbf{Frontend} --- aplikacja webowa zbudowana w \texttt{Next.js}, odpowiedzialna za interfejs użytkownika, wizualizację danych oraz interakcję z modułem backend.
\item \textbf{Backend} --- serwis napisany w Spring Boot, który obsługuje logikę biznesową, zarządza komunikacją z bazą danych oraz integruje się z modułem detekcji.
\item \textbf{Detector} --- moduł detekcji zachowań osób, napisany w Pythonie, który przetwarza strumienie wideo w czasie rzeczywistym, analizuje je pod kątem wykrywania określonych obiektów i zachowań oraz generuje zdarzenia do dalszego przetwarzania.
\item \textbf{PostgreSQL} --- relacyjna baza danych przechowująca informacje o zdarzeniach, konfiguracjach systemu oraz metadanych związanych z nagraniami wideo.
\end{itemize}